\documentclass[a4paper,11pt]{amsart}
\usepackage[margin=1in]{geometry}
\usepackage[T1]{fontenc}
\usepackage{lmodern}
\usepackage{amsmath,amssymb,amsthm}
\usepackage{microtype}
\usepackage[hidelinks]{hyperref}

\newtheorem{theorem}{Theorem}[section]
\newtheorem{proposition}[theorem]{Proposition}
\newtheorem{lemma}[theorem]{Lemma}
\newtheorem{corollary}[theorem]{Corollary}
\theoremstyle{remark}
\newtheorem{remark}[theorem]{Remark}
\DeclareMathOperator{\conv}{conv}
\newcommand{\R}{\mathbb R}
\newcommand{\C}{\mathbb C}
\newcommand{\ii}{\mathrm i}
\newcommand{\Hh}{\mathcal H}

\title{Unit distances in convex polygons}
\author{Liam Kruer}
\author{Jensen Kohlmeyer}
\author{Liam Price}
\date{}
\subjclass[2020]{Primary 52C10; Secondary 52A10, 05C62}
\keywords{Unit distances, convex position, minimum degree, Erd\H{o}s problems}

\begin{document}
\begin{abstract}
For every sufficiently large $n$, we construct $n$ points in strictly convex
position such that every point has at least
$(\frac14-o(1))\log_2\log_2 n$ unit-distance neighbours and the set determines
at least $(\frac14-o(1))n\log_2\log_2 n$ unordered unit-distance pairs.
The points can be confined to two arbitrarily small disks whose centres are
one unit apart, so their unit-distance graph is bipartite.
These results give negative answers to Erd\H{o}s Problems 96 and 97 and to
the more general conjecture that some universal bound on equal-distance
multiplicity holds at a suitably chosen vertex of every convex polygon.
The construction combines the middle levels graph, products of small
rotations with bounded exponents, and independent perturbations of the radii; vertex deletion
and unions of slightly rotated copies give the minimum-degree bounds at
every sufficiently large cardinality.
\end{abstract}
\maketitle

\section{Introduction}

A finite set $P\subset\R^2$ is in \emph{strictly convex position} if
$p\notin\conv(P\setminus\{p\})$ for every $p\in P$.
Thus every point is a vertex of its convex hull.
Write $u(P)$ for the number of unordered pairs of points of $P$ at distance
one, and, for nonempty $P$, put
\[
 d_1(p;P)=\#\{q\in P:\|p-q\|=1\},\qquad
 \delta_1(P)=\min_{p\in P}d_1(p;P).
\]
The \emph{unit-distance graph} of $P$ has vertex set $P$ and joins precisely
the pairs at distance one. In particular, $u(P)$ is its number of edges and
$\delta_1(P)$ is its minimum degree. Let $U_c(n)$ be the maximum of $u(P)$
over strictly convex $n$-point sets. All logarithms below are to base two,
and $B(c,\eta)$ denotes the open Euclidean disk with centre $c$ and radius
$\eta$.

Erd\H{o}s Problem 96 asks whether $U_c(n)=O(n)$.
Problem 97 asks whether every convex polygon has a vertex from which no
four other vertices are equidistant. The general version asks whether
there is an integer $k$ such that every convex polygon has a vertex from
which no $k$ other vertices are equidistant
\cite{Bloom96,Bloom97,FormalConjectures}.
The distance in these latter questions is allowed to depend on the vertex.
Our examples use the same distance, namely one, at every vertex.

\begin{theorem}\label{thm:main}
There are absolute constants $C>0$ and $n_0$ such that, for every integer
$n\ge n_0$ and every $0<\eta<1/2$, there is a strictly convex $n$-point set
$P\subset\R^2$ satisfying
\begin{equation}\label{eq:main}
 \min\left\{\delta_1(P),\frac{u(P)}n\right\}
 \ge \frac14\log\log n-C\log\log\log n.
\end{equation}
Moreover,
$P\subset B((0,-1/2),\eta)\cup B((0,1/2),\eta)$.
Its unit-distance graph is bipartite, and its two vertex classes are
separated by the horizontal axis.
\end{theorem}

\begin{corollary}\label{cor:problems}
We have $U_c(n)=\Omega(n\log\log n)$.
For every integer $k\ge1$ and every sufficiently large integer $n$, there
is a strictly convex $n$-point set in which every point has at least $k$
unit-distance neighbours. Both assertions hold with the two-disk
restriction in Theorem~\ref{thm:main}.
Consequently, Problem 96, Problem 97, and its general fixed-$k$ version
all have negative answers.
\end{corollary}

\begin{proof}
The right-hand side of \eqref{eq:main} tends to infinity and is
$(1/4-o(1))\log\log n$. The first assertion follows from the bound for
$u(P)$, and the second from the bound for $\delta_1(P)$.
Taking $k=4$ disproves Problem 97, while arbitrary $k$ disproves its
general version.
\end{proof}

Our construction starts from the middle levels graph. We realize its
incidences as exact unit distances after small rotations, and use
independent perturbations of the radii to make all bounded products of
the rotations distinct. A strict supporting-line inequality preserves
every point as a vertex of the convex hull. Vertex deletion and unions
of slightly rotated copies then give simultaneous minimum-degree and
edge-count bounds at every sufficiently large cardinality. The use of
commuting rotations and mixed derivatives to separate configurations is
adapted from Swanepoel and Valtr~\cite{SV}. For comparison,
Aggarwal~\cite{Aggarwal} proved $U_c(n)\le n\log n+4n$; our lower bound
does not determine the order of $U_c(n)$.

The following finite construction is the geometric input to the proof.
The extra supporting-line inequality will allow unions of slightly
rotated copies without losing any vertices.

\begin{theorem}\label{thm:finite}
For integers $d,q\ge2$, put
\begin{equation}\label{eq:parameters}
 v_d=\binom{2d}{d},\qquad m_d=\frac d2\binom{2d}{d},\qquad
 N_{d,q}=v_dq^{m_d}.
\end{equation}
For every $0<\eta<1/2$, there is a set $P\subset\R^2\setminus\{0\}$
with
\begin{equation}\label{eq:finite-count}
 |P|=N_{d,q},\qquad
 u(P)\ge \frac d2\left(1-\frac1q\right)N_{d,q},
\end{equation}
such that $P\subset B((0,-1/2),\eta)\cup B((0,1/2),\eta)$ and
\begin{equation}\label{eq:support}
 \langle x,y-x\rangle<0\qquad(x,y\in P,\ x\ne y).
\end{equation}
In particular, $P$ is in strictly convex position.
\end{theorem}

\section{Products of rotations}

Identify $\R^2$ with $\C$, with inner product
$\langle z,w\rangle=\operatorname{Re}(\overline z w)$, and define
$\Hh(z,w)=\langle z,w-z\rangle$.
If a finite set satisfies \eqref{eq:support}, the linear functional
$x\mapsto\langle z,x\rangle$ is uniquely maximized at $z$ for each point
$z$ of the set. This proves strict convexity.
The identity
$2\Hh(z,w)=|w|^2-|z|^2-|w-z|^2$ also shows that distinct points of equal
modulus satisfy $\Hh(z,w)=-|w-z|^2/2<0$.
For unit complex numbers $u,v$, we have
\begin{equation}\label{eq:rotation-bound}
 |\Hh(zu,wv)-\Hh(z,w)|
 \le |z||w|\bigl(|u-1|+|v-1|\bigr),
\end{equation}
because the difference equals
$\operatorname{Re}(\overline z w(\overline u v-1))$.

\begin{lemma}\label{lem:products}
Let $G=(L\sqcup R,E)$ be a finite simple bipartite graph with
$v=|L|+|R|$ and $m=|E|\ge1$; write the endpoints of $e$ as
$\ell(e)\in L$ and $r(e)\in R$.
Fix an integer $q\ge2$.
Suppose that nonzero complex numbers $b_i$, indexed by $L\sqcup R$,
have pairwise distinct moduli at most one, and that unit complex numbers
$w_e$, indexed by $E$, satisfy
\begin{equation}\label{eq:product-hypotheses}
 |b_{\ell(e)}-b_{r(e)}w_e|=1,\qquad
 \Hh(b_i,b_j)<-2H\quad(i\ne j),\qquad
 H=(q-1)\sum_{e\in E}|w_e-1|.
\end{equation}
Suppose also that the products
$W_{\mathbf a}=\prod_{e\in E}w_e^{a_e}$, where
$\mathbf a\in\{0,\ldots,q-1\}^{E}$, are pairwise distinct.
Then
\begin{equation}\label{eq:product-set}
 P=\{b_iW_{\mathbf a}:i\in L\sqcup R,\
                   \mathbf a\in\{0,\ldots,q-1\}^{E}\}
\end{equation}
satisfies \eqref{eq:support}, has $vq^m$ points, and determines at least
$m(q-1)q^{m-1}$ unit-distance pairs.
Moreover, $|b_iW_{\mathbf a}-b_i|\le H$ for every index.
\end{lemma}

\begin{proof}
Equality $b_iW_{\mathbf a}=b_jW_{\mathbf b}$ implies $i=j$ by taking
moduli, and then $\mathbf a=\mathbf b$ by cancellation and the assumed
distinctness of the products. Thus the indexing is injective.
Telescoping products gives
\[
 |W_{\mathbf a}-1|
 \le\sum_{e\in E}|w_e^{a_e}-1|
 \le\sum_{e\in E}a_e|w_e-1|\le H.
\]
For $i\ne j$, inequality \eqref{eq:rotation-bound} and
\eqref{eq:product-hypotheses} give
$\Hh(b_iW_{\mathbf a},b_jW_{\mathbf b})<0$.
For $i=j$ and $\mathbf a\ne\mathbf b$, the points are distinct and have
equal modulus, so the same conclusion follows from the preceding
identity for $\Hh$. The displacement bound follows from $|b_i|\le1$.

Let $\mathbf1_e$ be the coordinate vector at $e$.
For each $e\in E$ and each $\mathbf a$ with $a_e\le q-2$, the pair
\[
 \bigl\{b_{\ell(e)}W_{\mathbf a},\,
        b_{r(e)}W_{\mathbf a+\mathbf1_e}\bigr\}
\]
has distance one, since $W_{\mathbf a+\mathbf1_e}=W_{\mathbf a}w_e$.
There are $m(q-1)q^{m-1}$ such pairs. Equality of two of these unordered
pairs cannot interchange the label classes $L$ and $R$, by injectivity
of the vertex indexing. The two endpoint labels therefore agree, the
simplicity of $G$ identifies the edge, and the $L$ endpoint identifies
$\mathbf a$. Hence all the counted pairs are distinct.
\end{proof}

\section{The geometric construction}

Fix $d\ge2$ and let
$L=\binom{[2d-1]}{d-1}$ and $R=\binom{[2d-1]}d$, where
$[2d-1]=\{1,\ldots,2d-1\}$.
Join $A\in L$ to $B\in R$ when $A\subset B$.
This is the middle levels graph: it is the subgraph of the
$(2d-1)$-dimensional hypercube induced by its two middle levels.
Every lower-level vertex has $d$ extensions and every upper-level
vertex has $d$ immediate subsets. Thus the graph is simple and
$d$-regular, with $v_d$ vertices and $m_d$ edges as in
\eqref{eq:parameters}.
Set $z_j=5^{j-1}$ and, for every label $A\in L\sqcup R$, define
\begin{equation}\label{eq:digits}
 X_A=\sum_{j\in A}z_j,\qquad T_A=1+\sum_{j\in A}z_j^2.
\end{equation}
For an edge $A\subset B=A\cup\{j\}$ we have
$X_B-X_A=z_j>0$ and $T_B-T_A=(X_B-X_A)^2$.

\begin{lemma}\label{lem:digits}
Distinct labels have distinct $X$-values and distinct $T$-values, and
\begin{equation}\label{eq:digit-bound}
 |T_B-T_A|<\frac{50}{27}(X_B-X_A)^2\qquad(A\ne B).
\end{equation}
\end{lemma}

\begin{proof}
Let $h$ be the largest exponent occurring in the symmetric difference
of the two labels, and put $M=5^h$.
The dominant digit gives
\[
 |X_B-X_A|\ge M-\sum_{j=0}^{h-1}5^j>\frac34M,\qquad
 |T_B-T_A|\le M^2+\sum_{j=0}^{h-1}25^j<\frac{25}{24}M^2.
\]
These inequalities imply \eqref{eq:digit-bound} and distinctness of the
$X$-values. The lower bound
$|T_B-T_A|\ge M^2-\sum_{j=0}^{h-1}25^j>0$ proves distinctness of the
$T$-values.
\end{proof}

Write $C(x,t)=x^2+t/2$ and $D(x,t)=-x^2+t/2$.
For $\varepsilon>0$, define
\begin{equation}\label{eq:seed}
 \begin{aligned}
 a_A(\varepsilon)
   &=\varepsilon X_A+\ii\bigl(-\tfrac12+
                    \varepsilon^2 C(X_A,T_A)\bigr),&&A\in L,\\
 a_B(\varepsilon)
   &=-\varepsilon X_B+\ii\bigl(\tfrac12+
                    \varepsilon^2 D(X_B,T_B)\bigr),&&B\in R.
 \end{aligned}
\end{equation}
Throughout the construction, all asymptotic estimates refer to
$\varepsilon\to0^+$ with $d$ and the finite label sets fixed.
Their constants may depend on these data.

\begin{lemma}\label{lem:seed}
There are $c>0$ and $\varepsilon_0>0$ such that, for
$0<\varepsilon<\varepsilon_0$, the points in \eqref{eq:seed} have
pairwise distinct moduli in $(0,1)$ and satisfy
$\Hh(a_i(\varepsilon),a_j(\varepsilon))<-c\varepsilon^2$ for $i\ne j$.
For an edge with data $(x,t)=(X_A,T_A)$ and $(y,u)=(X_B,T_B)$,
\begin{equation}\label{eq:distance-error}
 |a_B(\varepsilon)-a_A(\varepsilon)|^2
 =1+\varepsilon^4\bigl(D(y,u)-C(x,t)\bigr)^2.
\end{equation}
\end{lemma}

\begin{proof}
For two labels of the same class, with data $(x,t)$ and $(y,u)$, direct
expansion gives, respectively,
\begin{align*}
 \Hh(a_A,a_{A'})
 &=\varepsilon^2\left(-\frac{(y-x)^2}{2}-\frac{u-t}{4}\right)
   +\varepsilon^4 C(x,t)\bigl(C(y,u)-C(x,t)\bigr),\\
 \Hh(a_B,a_{B'})
 &=\varepsilon^2\left(-\frac{(y-x)^2}{2}+\frac{u-t}{4}\right)
   +\varepsilon^4 D(x,t)\bigl(D(y,u)-D(x,t)\bigr).
\end{align*}
By \eqref{eq:digit-bound}, either quadratic coefficient is strictly less
than $-(y-x)^2/27<0$.
For labels from opposite classes, in either order, the support
expression tends to $-1/2$.
Since there are finitely many ordered pairs, these facts give a common
$c>0$ and a common sufficiently small interval with the claimed support
inequalities.
The exact modulus formulas are
\begin{equation}\label{eq:radii}
 \begin{aligned}
 |a_A|^2&=\tfrac14-\tfrac12\varepsilon^2T_A
              +\varepsilon^4 C(X_A,T_A)^2,&&A\in L,\\
 |a_B|^2&=\tfrac14+\tfrac12\varepsilon^2T_B
              +\varepsilon^4 D(X_B,T_B)^2,&&B\in R.
 \end{aligned}
\end{equation}
Within a class, distinct $T$-values give a nonzero quadratic coefficient
in the difference of any two squared moduli. Between classes that
coefficient is $(T_A+T_B)/2>0$. Thus all moduli are distinct for small
positive $\varepsilon$; they tend to $1/2$, so they lie in $(0,1)$.
Finally, if $M=D(y,u)-C(x,t)$, then
$|a_B-a_A|^2=\varepsilon^2(x+y)^2+(1+\varepsilon^2M)^2$.
On an edge, $u-t=(y-x)^2$, so $(x+y)^2+2M=0$, proving
\eqref{eq:distance-error}.
\end{proof}

\begin{lemma}\label{lem:correction}
For every edge $e=(A,B)$, and all sufficiently small positive
$\varepsilon$, there is a continuous choice of a unit complex number
$w_e(\varepsilon)$ such that
\begin{equation}\label{eq:correction}
 |a_A-a_Bw_e|=1,\qquad
 w_e-1=O(\varepsilon^3),\qquad
 \operatorname{Im}(\overline{a_A}a_Bw_e)<0.
\end{equation}
\end{lemma}

\begin{proof}
Fix an edge, abbreviate its data by $(x,t),(y,u)$, and put
$C=C(x,t)$, $D=D(y,u)$, and $M=D-C$.
Write
$Z=\overline{a_A}a_B=\xi+\ii\varepsilon\tau$, where
$\tau=(x-y)/2+\varepsilon^2(xD+yC)$.
Thus $Z\to-1/4$ and $\tau\to(x-y)/2<0$.
Define
\begin{equation}\label{eq:exact-w}
 F=\tau^2-\varepsilon^2M^2\xi-\tfrac14\varepsilon^6M^4,
 \qquad
 w=\frac{\xi+\tfrac12\varepsilon^4M^2-
                  \ii\varepsilon\sqrt F}{Z}.
\end{equation}
For small positive $\varepsilon$, both $Z\ne0$ and $F>0$ hold.
The identity
$(\xi+\varepsilon^4M^2/2)^2+\varepsilon^2F
 =\xi^2+\varepsilon^2\tau^2=|Z|^2$
shows that $|w|=1$. Moreover,
\[
 |a_A-a_Bw|^2
 =|a_A-a_B|^2-\varepsilon^4M^2=1,\qquad
 \operatorname{Im}(Zw)=-\varepsilon\sqrt F<0.
\]
To bound the correction, observe that $\sqrt F-\tau$ has a positive
limit and
$\sqrt F+\tau=(F-\tau^2)/(\sqrt F-\tau)=O(\varepsilon^2)$.
Consequently
$w-1=Z^{-1}(\varepsilon^4M^2/2-
\ii\varepsilon(\sqrt F+\tau))=O(\varepsilon^3)$.
All defining functions are continuous on a sufficiently small positive
interval.
\end{proof}

We next separate the bounded products while preserving the strict
inequalities and the exact unit distances.
For $r,s>0$ satisfying $-1<c(r,s)<1$, define
\begin{equation}\label{eq:angle}
 c(r,s)=\frac{r^2+s^2-1}{2rs},\qquad
 \gamma(r,s)=\arccos c(r,s)\in(0,\pi).
\end{equation}
This is the angle opposite the unit side of a nondegenerate triangle
with other side lengths $r,s$.

\begin{lemma}\label{lem:separation}
Let $G=(L\sqcup R,E)$ be a finite simple bipartite graph, and let
$U\subset(0,\infty)^{L\sqcup R}$ be nonempty and open, such that
$-1<c(\rho_{\ell(e)},\rho_{r(e)})<1$ for every edge throughout $U$.
Fix real numbers $\theta_e$ and an integer $q\ge2$, and put
\[
 w_e(\boldsymbol\rho)
 =\exp\bigl(\ii(\theta_e-
       \gamma(\rho_{\ell(e)},\rho_{r(e)}))\bigr).
\]
There is a vector in $U$ for which the products
$\prod_{e\in E}w_e(\boldsymbol\rho)^{a_e}$,
$\mathbf a\in\{0,\ldots,q-1\}^{E}$, are pairwise distinct.
\end{lemma}

\begin{proof}
First, throughout the domain in \eqref{eq:angle},
\begin{equation}\label{eq:mixed}
 \gamma_{rs}(r,s)
 =-\frac{c_{rs}(1-c^2)+c\,c_r c_s}{(1-c^2)^{3/2}}
 =\frac{1}{r^2s^2(1-c^2)^{3/2}}>0.
\end{equation}
Indeed,
$c_r=(r^2-s^2+1)/(2r^2s)$,
$c_s=(s^2-r^2+1)/(2rs^2)$, and
$c_{rs}=-(r^2+s^2+1)/(2r^2s^2)$.
Substituting these expressions gives
$-c_{rs}(1-c^2)-c\,c_r c_s=1/(r^2s^2)$, proving
\eqref{eq:mixed} by differentiation of $\arccos c$.

For distinct exponent vectors $\mathbf a,\mathbf b$, let
$\lambda_e=a_e-b_e$, and define
\[
 \Phi(\boldsymbol\rho)
 =\sum_{e\in E}\lambda_e\theta_e
  -\sum_{e\in E}\lambda_e
       \gamma(\rho_{\ell(e)},\rho_{r(e)}).
\]
The two products agree exactly on the relatively closed subset of $U$
where $\exp(\ii\Phi)=1$.
This subset has empty interior. Otherwise there would be a nonempty
open ball on which the continuous function $\Phi$ takes values in
$2\pi\mathbb Z$, and hence is constant. Choose $e_0$ with
$\lambda_{e_0}\ne0$. Because the graph is simple, every summand other
than $e_0$ depends on at most one of the two coordinates
$\rho_{\ell(e_0)},\rho_{r(e_0)}$.
The corresponding mixed derivative of $\Phi$ is therefore
$-\lambda_{e_0}\gamma_{rs}
 (\rho_{\ell(e_0)},\rho_{r(e_0)})$, which is nonzero by
\eqref{eq:mixed}. This contradicts constancy.
There are only finitely many pairs of exponent vectors.
A finite union of relatively closed sets with empty interior cannot
cover a nonempty open set: remove the sets successively, retaining a
nonempty relatively open subset at each step.
A point outside their union has the required property.
\end{proof}

\begin{proof}[Proof of Theorem~\ref{thm:finite}]
Fix $d,q$ and $\eta$, and construct the graph, seed points, and
corrections above. Put $c_i=-\ii/2$ for $i\in L$ and $c_i=\ii/2$ for
$i\in R$.
By Lemmas~\ref{lem:seed} and \ref{lem:correction},
$H_0=(q-1)\sum_e|w_e(\varepsilon)-1|=O(\varepsilon^3)$,
whereas the common support margin has order $\varepsilon^2$.
Thus one can fix a sufficiently small positive $\varepsilon$ for which
the moduli are pairwise distinct and in $(0,1)$, and
\begin{equation}\label{eq:open-conditions}
 \Hh(a_i,a_j)<-2H_0\quad(i\ne j),\qquad
 |a_i-c_i|+H_0<\eta\quad\text{for all }i.
\end{equation}
The second condition follows also from $a_i\to c_i$.
This choice is made after fixing $d$ and $q$; no estimate uniform in
those parameters is required.

Write $a_i=\rho_i^0\alpha_i$ with $|\alpha_i|=1$ and
$\rho_i^0>0$, and keep the directions $\alpha_i$ fixed.
For $e=(A,B)$ choose a real $\theta_e$ satisfying
$\exp(\ii\theta_e)=\alpha_A/\alpha_B$.
At the original radii, \eqref{eq:correction} gives
\[
 \operatorname{Re}(\overline{a_A}a_Bw_e)
   =\frac{(\rho_A^0)^2+(\rho_B^0)^2-1}{2},\qquad
 |\overline{a_A}a_Bw_e|=\rho_A^0\rho_B^0.
\]
The strictly negative imaginary part gives
$-1<c(\rho_A^0,\rho_B^0)<1$. Writing
$\gamma_e^0=\gamma(\rho_A^0,\rho_B^0)$, we have
$\overline{a_A}a_Bw_e=\rho_A^0\rho_B^0e^{-\ii\gamma_e^0}$;
hence $w_e=e^{\ii(\theta_e-\gamma_e^0)}$ at the original vector.

For nearby positive radii $\boldsymbol\rho$, set
$b_i=\rho_i\alpha_i$ and use the functions $w_e(\boldsymbol\rho)$ in
Lemma~\ref{lem:separation}. Whenever all edge pairs are feasible,
\[
 |b_A-b_Bw_e(\boldsymbol\rho)|^2
 =\rho_A^2+\rho_B^2-
                  2\rho_A\rho_B\cos\gamma(\rho_A,\rho_B)=1.
\]
Write $H(\boldsymbol\rho)=(q-1)\sum_e|w_e(\boldsymbol\rho)-1|$.
All other required conditions are strict and continuous at
$\boldsymbol\rho^0$.
There is therefore a nonempty open neighbourhood $U$ on which the
radii remain pairwise distinct and in $(0,1)$, every edge pair is
feasible, and
\[
 \Hh(b_i,b_j)<-2H(\boldsymbol\rho)\quad(i\ne j),\qquad
 |b_i-c_i|+H(\boldsymbol\rho)<\eta.
\]
Lemma~\ref{lem:separation} supplies a vector in $U$ with all bounded
products distinct. Applying Lemma~\ref{lem:products} gives a set $P$
satisfying \eqref{eq:support}, with $v_dq^{m_d}$ points and at least
$m_d(q-1)q^{m_d-1}$ unit-distance pairs. Since $m_d/v_d=d/2$,
these are exactly the bounds in \eqref{eq:finite-count}.
The displacement bound in that lemma gives
$|b_iW_{\mathbf a}-c_i|\le H(\boldsymbol\rho)+|b_i-c_i|<\eta$,
which proves the required localization.
\end{proof}

\section{Minimum degree and prescribed cardinality}

\begin{lemma}\label{lem:deletion}
Let $G$ be a finite nonempty graph with
$|E(G)|\ge\alpha|V(G)|$, where $\alpha>k$ and $k$ is a nonnegative
integer. Then $G$ has a nonempty induced subgraph $H$ with minimum
degree at least $k+1$ and $|E(H)|\ge\alpha|V(H)|$.
\end{lemma}

\begin{proof}
Repeatedly delete a vertex of current degree at most $k$ whenever one
exists. Deleting a vertex of degree $t\le k$ increases the quantity
$|E|-\alpha|V|$ by $\alpha-t>0$.
The process cannot delete every vertex: that quantity is initially
nonnegative, would become positive after the first deletion, but is
zero for the empty graph.
The surviving graph has the asserted minimum degree, and the
nonnegativity of the same quantity gives its edge-count bound.
\end{proof}

\begin{lemma}\label{lem:copies}
Let $Q\subset\C\setminus\{0\}$ be finite and satisfy
\eqref{eq:support}, and let $\mathcal O\subset\C$ be open with
$Q\subset\mathcal O$.
For any finite list of subsets $Q_1,\ldots,Q_t\subseteq Q$, there are
rotations about the origin whose images of these subsets are pairwise
disjoint and whose union lies in $\mathcal O$ and satisfies
\eqref{eq:support}.
The rotation angles may be chosen arbitrarily close to zero.
\end{lemma}

\begin{proof}
Finiteness, continuity, and the strict support inequalities give a
number $0<\varepsilon<\pi$ such that, whenever $|\theta|,|\phi|<\varepsilon$,
we have $e^{\ii\theta}x\in\mathcal O$ for every $x\in Q$ and
$\Hh(e^{\ii\theta}x,e^{\ii\phi}y)<0$ for all distinct $x,y\in Q$.
Choose distinct angles $\theta_1,\ldots,\theta_t$ in this interval.
Points arising from different elements of $Q$ satisfy the strict
support inequality, so cannot coincide.
For a fixed $x\in Q$ and $j\ne\ell$, the points
$e^{\ii\theta_j}x$ and $e^{\ii\theta_\ell}x$ are distinct because
$x\ne0$ and $0<|\theta_j-\theta_\ell|<2\pi$.
They have equal modulus, so they too satisfy the strict support
inequality in both orders. This proves all the assertions.
\end{proof}

\begin{proposition}\label{prop:threshold}
For $k\ge1$, put $M_k=N_{2k+1,\,2k+2}$.
For every integer $n\ge M_k^2$ and every $0<\eta<1/2$, there is a set
$P$ satisfying \eqref{eq:support}, contained in
$B((0,-1/2),\eta)\cup B((0,1/2),\eta)$, such that
\begin{equation}\label{eq:threshold}
 |P|=n,\qquad \delta_1(P)\ge k,\qquad u(P)\ge(k-1)n.
\end{equation}
\end{proposition}

\begin{proof}
Apply Theorem~\ref{thm:finite} with $d=2k+1$ and $q=2k+2$.
Its set has $M_k$ points and at least $\alpha M_k$ unit-distance pairs,
where
\[
 \alpha=\frac{2k+1}{2}\left(1-\frac1{2k+2}\right)
        =k+\frac1{4k+4}>k.
\]
Apply Lemma~\ref{lem:deletion} to its full unit-distance graph.
The surviving vertex set $Q$ has size $s\le M_k$, minimum degree at
least $k+1$, and $u(Q)\ge\alpha s$.
It inherits \eqref{eq:support} and the two-disk restriction.
In particular $s\ge k+2$.
Choose a vertex $v$ of $Q$ whose degree is at most the average degree
$2u(Q)/s$, and let $Q^-=Q\setminus\{v\}$.
Then $|Q^-|=s-1$ and $\delta_1(Q^-)\ge k$. Moreover,
\[
 \frac{u(Q^-)}{s-1}
 \ge\frac{u(Q)(s-2)}{s(s-1)}
 \ge\alpha\frac{s-2}{s-1}
 >k-\frac{k}{s-1}>k-1.
\]
Thus both $Q$ and $Q^-$ have minimum degree at least $k$ and at least
$k-1$ unit-distance edges per vertex.

Since $n\ge M_k^2\ge s^2$, there are nonnegative integers $a,b$ with
$n=as+b(s-1)$.
Indeed, choose $b\in\{0,\ldots,s-1\}$ with $b\equiv-n\pmod s$;
then $a=(n-b(s-1))/s$ is an integer and is nonnegative because
$n\ge s^2>(s-1)^2\ge b(s-1)$.
Use Lemma~\ref{lem:copies}, with $\mathcal O$ equal to the union of the
two prescribed disks, on $a$ copies of $Q$ and $b$ copies of $Q^-$.
Their disjoint rotated images form a set $P$ with exactly $n$ points
satisfying \eqref{eq:support} and the required localization.
Every vertex retains its unit-distance neighbours within its own
copy, so $\delta_1(P)\ge k$.
Counting just the edges within individual copies gives
$u(P)\ge a(k-1)s+b(k-1)(s-1)=(k-1)n$.
Any additional edges between copies can only increase these bounds.
\end{proof}

\begin{proof}[Proof of Theorem~\ref{thm:main}]
The elementary binomial estimates
$4^d/(2d+1)\le\binom{2d}{d}\le4^d$ follow by comparing the largest
coefficient in $(1+1)^{2d}$ with the sum of all its coefficients.
Together with \eqref{eq:parameters}, they give
\[
 \log N_{d,d+1}
 =\log v_d+\frac d2v_d\log(d+1),\qquad
 \log\log\bigl(N_{d,d+1}^{\,2}\bigr)=2d+O(\log d).
\]
Consequently, for the strictly increasing sequence $M_k$ in
Proposition~\ref{prop:threshold},
\begin{equation}\label{eq:threshold-growth}
 T_k:=\log\log(M_k^2)=4k+O(\log(k+1)).
\end{equation}
All constants here are absolute.
For sufficiently large $n$, let $k$ be the largest integer with
$M_k^2\le n$, and write $L=\log\log n$.
Then $T_k\le L<T_{k+1}$.
The lower estimate in \eqref{eq:threshold-growth} first gives
$k=O(L)$, and the upper estimate then gives
$k\ge L/4-O(\log L)$.
Proposition~\ref{prop:threshold} therefore supplies a set with
\[
 \min\left\{\delta_1(P),\frac{u(P)}n\right\}
 \ge k-1\ge\frac14\log\log n-C\log\log\log n
\]
for a suitable absolute $C$ and all $n$ beyond an absolute threshold.
The construction works for every prescribed $\eta$ without changing
that threshold.
Each of the two disks has diameter $2\eta<1$, so it contains no
unit-distance pair. Both disks lie strictly on opposite sides of the
horizontal axis. Their intersections with $P$ thus give the claimed
bipartition and separation.
\end{proof}

\begin{corollary}\label{cor:size}
For every integer $k\ge2$ and every $0<\eta<1/2$, there is a strictly
convex set $Q$ in the two prescribed disks with
$\delta_1(Q)\ge k$ and
\begin{equation}\label{eq:size}
 |Q|\le N_{2k-1,\,2k}=2^{\,2^{\,4k+O(\log k)}}.
\end{equation}
In particular, a counterexample to Problem 97 exists with at most
$3432\cdot2^{36036}$ vertices.
\end{corollary}

\begin{proof}
Use Theorem~\ref{thm:finite} with $d=2k-1$ and $q=2k$.
The number of unit-distance pairs per vertex is at least
$\frac{2k-1}{2}(1-\frac1{2k})=k-1+\frac1{4k}>k-1$.
Lemma~\ref{lem:deletion}, applied with deletion threshold $k-1$, leaves
a nonempty induced subgraph of minimum degree at least $k$.
Its vertex set inherits the geometric properties and has at most
$N_{2k-1,2k}$ points. Since $k\ge2$, it has at least three vertices.
The growth estimate follows from the same binomial bounds used above.
For $k=4$, the parameters are $d=7$, $q=8$,
$v_7=\binom{14}{7}=3432$, and $m_7=12012$; hence
$N_{7,8}=3432\cdot8^{12012}=3432\cdot2^{36036}$.
\end{proof}

\begin{remark}
The existence of strictly convex sets with arbitrarily large minimum
unit-distance degree is equivalent to the unboundedness of
$u(P)/|P|$ over strictly convex sets.
One direction follows from $2u(P)=\sum_{p\in P}d_1(p;P)$;
the other follows from Lemma~\ref{lem:deletion}, since strict convex
position is inherited by subsets.
Thus the existence-only consequences for the two Erd\H{o}s problems are
logically linked. The quantitative simultaneous bounds and the control
of every sufficiently large cardinality are the additional conclusions
of Theorem~\ref{thm:main}.
The perturbation used in Theorem~\ref{thm:finite} is an existence
argument avoiding finitely many closed exceptional sets, not a
numerical approximation. All counted distances are exactly one.
The explicit quantities in \eqref{eq:parameters} and \eqref{eq:size}
are cardinality bounds, not bounds on the bit complexity of coordinates.
\end{remark}

\section*{AI disclosure}
The authors report that the argument originated in a disproof of
Erd\H{o}s Problem~96 generated entirely by GPT~6 Astra, including the
construction and its mathematical justification. The quantitative
strengthening, the consequences for Problem~97 and its general
fixed-$k$ version, and the extension to every sufficiently large
cardinality were subsequently developed through further interaction
with the model. The model also produced the exposition and
\LaTeX{} manuscript. This disclosure describes the provenance of the
argument; it does not supersede the attribution of previously
published mathematical methods cited in the paper or constitute
independent verification of the proofs. The human authors retain
responsibility for the correctness of the mathematical claims, the
accuracy of this disclosure, and the decision to submit or
disseminate the work.

\begin{thebibliography}{9}

\bibitem{Aggarwal}
A. Aggarwal,
\emph{On unit distances in a convex polygon},
Discrete Math. \textbf{338} (2015), 88--92.

\bibitem{Bloom96}
T. F. Bloom,
\emph{Erd\H{o}s Problem \#96},
\url{https://www.erdosproblems.com/96}.

\bibitem{Bloom97}
T. F. Bloom,
\emph{Erd\H{o}s Problem \#97},
\url{https://www.erdosproblems.com/97}.

\bibitem{FormalConjectures}
Google DeepMind and the Formal Conjectures contributors,
\emph{Formal Conjectures},
entries for Erd\H{o}s Problems~96 and~97,
\url{https://github.com/google-deepmind/formal-conjectures},
accessed September 13, 2026.

\bibitem{SV}
K. J. Swanepoel and P. Valtr,
\emph{The unit distance problem on spheres},
in \emph{Towards a Theory of Geometric Graphs} (J. Pach, ed.),
Contemp. Math. \textbf{342}, Amer. Math. Soc., Providence, RI, 2004,
273--279.

\end{thebibliography}
\end{document}